\subsection{作用的逆向极限}

假设给定相对于同一指标集 I 的两个集合逆向系统 \((\Omega_{\alpha}, \phi_{\alpha\beta})\) 与 \((\mathrm{E}_{\alpha}, f_{\alpha\beta})\) 。假设对所有 \(\alpha \in I\) 给定了 \(\Omega_{\alpha}\) 在 \(E_{\alpha}\) 上的一个作用，使得

\[
f _ {\alpha \beta} (\omega_ {\beta} x _ {\beta}) = \phi_ {\alpha \beta} (\omega_ {\beta}) \cdot f _ {\alpha \beta} (x _ {\beta})\tag{1}
\]

对 \(\alpha \leqslant \beta, x_{\beta} \in \mathrm{E}_{\beta}, \omega_{\beta} \in \Omega_{\beta}\) 。那么所考虑的这族作用称为一个作用逆向系统。设 \(\Omega = \varprojlim \Omega_{\alpha}\) 与 \(\mathrm{E} = \varprojlim \mathrm{E}_{\alpha}\) ；若 \(x = (x_{\alpha})_{\alpha \in I}\) 属于 \(\mathrm{E}\) 且 \(\omega = (\omega_{\alpha})_{\alpha \in I}\) 属于 \(\Omega\) ，则 \(\omega \cdot x = (\omega_{\alpha} \cdot x_{\alpha})_{\alpha \in I}\) 由(1)属于 E。这样就定义了 \(\Omega\) 在 E 上的一个作用，称为诸 \(\Omega_{\alpha}\) 在诸 \(\mathrm{E}_{\alpha}\) 上作用的逆极限。

上述内容特别适用于 \(\Omega_{\alpha}\) 是幺半群且每个 \(\Omega_{\alpha}\) 在 \(E_{\alpha}\) 上的作用是一个运算的情形。那么这些运算的逆极限是该幺半群在 E 上的一个运算。

留给读者定义带算子的群的逆系统的逆向极限，并验证该极限是带算子的群。

